\begin{abstract}
On-Policy Self-Distillation (OPSD) enables a model to learn from a privileged version of its
own context without requiring a separate larger teacher model. While OPSD has primarily been
investigated for language-model reasoning, its effectiveness for vision-language models for visual question-answering remains
less understood. This work evaluates OPSD for visual question answering using the GQA dataset. Question-grounded
object annotations are used to construct privileged image crops for the teacher while the student
receives the full image. OPSD is compared with conventional Supervised Fine-Tuning (SFT) across
SmolVLM-500M-Instruct, BLIP-VQA-Base and MiniCPM-V 4.6. Results show that SFT substantially outperformed
OPSD for SmolVLM and BLIP, while both methods achieved similar BERTScore-F1 performance for MiniCPM.
The crop-conditioned teacher also provides only a weak average advantage over the full-image model, limiting
the strength of the distillation signal. Qualitative analysis shows that privileged crops can help when
relevant visual evidence is localised within a complex image, but may remove context required for
broader scene understanding. These findings suggest that OPSD performance is strongly dependent on both
model characteristics and the quality of the privileged context. While SFT remains the stronger general-purpose
adaptation strategy in this setting, context-based self-distillation may be more effective in domains
where privileged training-time information provides a clearer and more reliable advantage.
\end{abstract}